\chapter{Frostbite}

\section*{Definition and Overview}
Frostbite is injury to the skin and underlying tissues caused by freezing, most often affecting the fingers, toes, ears, nose, and cheeks. It occurs when skin is exposed to temperatures below freezing, especially with wind, and the tissues freeze because blood flow is constricted to preserve core warmth. Frostbite ranges from mild frostnip, which recovers fully, to deep frostbite with tissue death (gangrene) that can require amputation. It is most common in cold climates, at high altitude, and in people who are homeless, intoxicated, or involved in winter sport, and it is a risk in parts of Australia's alpine regions. This chapter explains how frostbite happens, how to recognise it, and how it is treated.

\section*{Epidemiology}
Frostbite is uncommon in most of Australia but occurs in alpine areas and among travellers to cold climates, high-altitude mountaineers, and people in cold outdoor work. Globally, it is most common in people who are homeless, people affected by alcohol or drugs, and outdoor workers and athletes. Risk is greatly increased by wind chill, wet clothing, fatigue, and pre-existing conditions that impair circulation, including diabetes, smoking, and peripheral vascular disease. While deep frostbite with amputation is the feared outcome, most cases are superficial and heal fully with prompt rewarming.

\section*{Aetiology and Pathophysiology}
\begin{itemize}
    \item \textbf{Freezing:} exposure of skin to temperatures below about minus 0.5 degrees, especially with wind chill and moisture, freezes the tissue water
    \item \textbf{Vascular response:} the body constricts blood vessels in the extremities to preserve heat for the core, reducing blood flow and worsening local injury
    \item \textbf{Ice crystal formation:} ice crystals form inside and between cells, damaging cell membranes and dehydrating cells
    \item \textbf{Blood flow damage:} on rewarming, blood returns to damaged vessels, causing inflammation, swelling, and clot formation that can destroy tissue
    \item \textbf{Depth of injury:} superficial frostbite affects the skin; deep frostbite involves deeper tissues, including bone, with tissue death
    \item \textbf{Risk factors:} cold temperature, wind chill, wet or tight clothing, alcohol, smoking, fatigue, dehydration, and poor circulation
    \item \textbf{The danger of refreezing:} tissue that thaws then refreezes is much more severely damaged
    \item \textbf{Associated hypothermia:} frostbite often occurs with a dangerously low core body temperature, which is the greater threat to life
\end{itemize}

\section*{Clinical Features}
\begin{itemize}
    \item Early (frostnip): cold, numb, pale skin, which tingles and reddens on rewarming, with no lasting damage
    \item Superficial frostbite: white or waxy skin that is numb, with a firm surface but soft deeper tissue; blisters form within 24 hours
    \item Deep frostbite: hard, waxy skin that does not move over the joints, with no sensation; large blood-filled blisters, then blackened tissue (gangrene)
    \item Affected parts: fingers, toes, ears, nose, cheeks, and chin
    \item Pain, which may be absent early because of numbness, then severe burning and throbbing on rewarming
    \item Swelling and discolouration after rewarming
    \item Loss of function of the affected part
    \item Signs of hypothermia: shivering, confusion, drowsiness, and slurred speech
\end{itemize}

\section*{Differential Diagnosis}
Cold injury has other forms that must be distinguished.
\begin{itemize}
    \item Frostnip, a mild, fully reversible form
    \item Hypothermia, low core body temperature, which can occur with or without frostbite
    \item Chilblains, painful red or purple swellings from cold without freezing
    \item Trench foot, injury from prolonged cold, wet exposure without freezing
    \item Raynaud's phenomenon and other circulatory conditions
    \item The depth of frostbite is assessed by the appearance and by its response to rewarming
\end{itemize}

\section*{When to Seek Medical Care}
Frostbite should be assessed by a doctor, especially if the skin is hard, white, and numb, or if it involves a significant area. Any frostbite with suspected hypothermia is an emergency.

\textbf{Call 000 (or local emergency services) for:}
\begin{itemize}
    \item Signs of hypothermia: shivering, confusion, drowsiness, or loss of consciousness
    \item Deep frostbite with hard, waxy, numb skin or blackened areas
    \item Frostbite of a large area or in a child or older person
    \item Any emergency in a cold-exposure setting
\end{itemize}

\section*{Diagnosis and Investigations}
\begin{itemize}
    \item \textbf{History and examination:} the exposure, duration, and appearance of the affected parts, including their colour, temperature, sensation, and capillary return
    \item \textbf{Assessment after rewarming:} the full extent of damage is often only clear days to weeks later
    \item \textbf{Imaging:} bone scans and angiography in the days after injury identify which tissue is viable and guide amputation decisions
    \item \textbf{Assessment for hypothermia:} core temperature measurement and treatment
    \item \textbf{Blood tests:} for dehydration, muscle breakdown, and infection
\end{itemize}

\section*{Management}
\begin{itemize}
    \item \textbf{Rewarming:} immerse the affected part in warm (not hot) water at 37--39 degrees for 15--30 minutes, until it is red and pliable; do not rub the area or use direct heat
    \item \textbf{Do not thaw if there is a risk of refreezing:} tissue that thaws then refreezes is more severely damaged
    \item \textbf{Pain relief:} rewarming is painful and requires analgesics
    \item \textbf{Blisters:} left intact by medical staff; large blood-filled blisters may be drained under sterile conditions
    \item \textbf{Medicines:} anti-inflammatory drugs and treatments to improve blood flow; antibiotics for infection
    \item \textbf{Tissue saving:} treatment in a specialist unit, with decisions about amputation delayed until the full extent of damage is clear, often weeks later
    \item \textbf{Hydrotherapy and physiotherapy:} to maintain movement during recovery
    \item \textbf{Treat hypothermia:} gradual rewarming of the whole body, which takes priority
    \item \textbf{Tetanus vaccination and wound care} as needed
\end{itemize}

\section*{Complications}
\begin{itemize}
    \item Permanent numbness, cold sensitivity, and pain in the affected parts
    \item Gangrene and amputation of fingers, toes, or larger parts of limbs in deep frostbite
    \item Infection of damaged tissue
    \item Chronic pain and altered sensation lasting years
    \item The complications of hypothermia, including heart rhythm problems
\end{itemize}

\section*{Prognosis and Outcomes}
The outcome depends on the depth and duration of the freezing. Frostnip and superficial frostbite heal fully, though the affected areas may remain sensitive to cold for months or years. Deep frostbite can cause lasting damage and amputation, but modern management, including delayed surgical decisions and treatments that restore blood flow, saves many tissues that would previously have been lost. With prompt rewarming, correct treatment, and rehabilitation, the majority of people recover good function of the affected parts.

\section*{Prevention and Self-Care}
\begin{itemize}
    \item Dress for cold: layers, insulated gloves, warm footwear, and a hat that covers the ears
    \item Keep clothing dry, as wet clothing loses insulation
    \item Check exposed skin regularly and come in from the cold at the first sign of numbness
    \item Avoid alcohol before and during cold exposure
    \item Stay hydrated, well fed, and rested
    \item Do not smoke, as it impairs circulation
    \item If frostbite occurs, rewarm gently in warm water and seek medical care
    \item Never rub frozen skin or use direct heat
\end{itemize}

\section*{Sources and Further Reading}
The following reputable sources provide further information for healthcare students and patients seeking reliable, current advice about frostbite.
\begin{itemize}
    \item Healthdirect Australia. \textit{Frostbite.} https://www.healthdirect.gov.au/frostbite
    \item Australian Red Cross. \textit{First aid for frostbite.} https://www.redcross.org.au/
    \item St John Ambulance Australia. \textit{Cold injuries.} https://stjohn.org.au/
    \item NHS. \textit{Frostbite.} https://www.nhs.uk/conditions/frostbite/
    \item Mayo Clinic. \textit{Frostbite.} https://www.mayoclinic.org/
    \item MedlinePlus. \textit{Frostbite.} https://medlineplus.gov/
\end{itemize}
